\def\Re{{\rm Re}}
\subsubsection{Infrared subtraction schemes}\seclab{IRsubtraction}

When plotting the elastic TPE contribution to $\delta_{2\gamma}$ or $R_{2\gamma}$, defined in Eqs.~\ref{delta2gTPE} and \ref{eq:R2gammadef}, one has to somehow treat the IR divergences they contain.
Of course the only way that leads to a measurable quantity is to combine them with the corresponding real corrections as discussed in Sec.~\ref{Sec:RadCorrAdrian}.
However, this can be quite cumbersome in practise and makes comparing different TPE scenarios and parametrisations needlessly complicated.
Instead, one defines an IR subtraction scheme that removes the IR divergence (and potentially some finite parts as well).
The exact form of this subtraction term is somewhat arbitrary and solely determined by conventions since the result is no longer a physical quantity.

A common way of doing this in the high-energy community is the definition of the Catani $I$ operator~\cite{Catani:1998bh} that simply removes the $1/\epsilon$ poles in dimensional regularisation.

In the QED community it is not uncommon to add the first terms of Eqs.~\ref{slicing} or \ref{subtraction}
\begin{eqnarray}
    \delta_{2\gamma}^{\rm eik} &= \delta_{2\gamma} +&
          \int_{E_\gamma < \delta} \dd\Phi_\gamma\ \mathcal{E} \nn\\
    & = \delta_{2\gamma} +&
          \underbrace{
              \int \dd\Phi_\gamma\Big(\frac{p_1\cdot k_1}{p_1\cdot k\ k_1\cdot k}+\frac{p_2\cdot k_2}{p_2\cdot k\ k_2\cdot k}}_{\hat{\mathcal{E}}(s)}\nn\\
              &&
         \hspace*{1cm} - (s\rightarrow u)\Big).
\end{eqnarray}
This scheme, sometimes called eikonal subtraction~\cite{Engel:2019nfw}, has the added advantage of giving the remnant a physical interpretation as long as the parameter $\delta$ is chosen small enough -- it just approximately includes real corrections up to the cut-off $\delta$.
The calculation of this integral is somewhat involved but the result is well known~\cite{Engel:2019nfw,Frederix:2009yq,Ulrich:2020phd}.
Since this physical interpretation exists, it also means that the result needs to be independent of how the TPE calculation was regularised, making comparisons between different methods easier.
The TPE community usually follows either Mo-Tsai (MoT) \cite{Mo:1968cg} or Maximon-Tjon (MTj) \cite{Maximon:2000hm}.
Note that in this paper, we use the MTj convention for our Figs.~\ref{fig.Dsig}-\ref{fig.NearForw}. The latter defines the IR-subtracted TPE contribution as
\begin{align}    &\delta_{2\gamma}^{\rm MTj} = \delta_{2\gamma} -
        \frac{2 \alpha}{\pi} q^2 \,\Re\Bigg[p_1\cdot k_1 \nn\\
        &\times\int  \frac{\dd \ell}{
            \ell^2
            (\ell-q)^2
            \Big[(\ell-k_1)^2-m^2\Big]
            \Big[(\ell+p_1)^2-M^2\Big]
        } \nn\\&\qquad\quad - (s\rightarrow u)\Bigg].
\end{align}
This integral can be evaluated both using photon-mass regularisation or dimensional regularisation
    \begin{eqnarray}
    \delta_{2\gamma}^{\rm MTj} &=& \delta_{2\gamma} +
    \frac{2 \alpha}{\pi}\,\Re\Bigg[\nn\\
        \frac{m^2+M^2-s}{\sqrt{\lambda_s}}
        & \times&\ln\frac{m^2+M^2-s + \sqrt{\lambda_s}}{2mM}
    \begin{cases} \ln\frac{m_\gamma^2}{Q^2} \\ \frac1{\epsilon}\Big(\frac{\mu^2}{Q^2}\Big)^\epsilon \end{cases}\nn\\
    &
    &- (s\rightarrow u)\Bigg],
\end{eqnarray}
where we have defined $\lambda_s = m^4 + (M^2 - s)^2 - 2 m^2 (M^2 + s)$.
In the limit of $m^2\ll s,M^2,Q^2$, the expression for $e^- p$ scattering simplifies to
    \begin{equation}
    \delta_{2\gamma}^{\rm MTj} = \delta_{2\gamma} +
    \frac{2 \alpha}{\pi}
        \ln\Big|\frac{s - M^2}{u - M^2}\Big|
    \begin{cases} \ln\frac{m_\gamma^2}{Q^2} \\ \frac1{\epsilon}\Big(\frac{\mu^2}{Q^2}\Big)^\epsilon \end{cases}\!\!\!\!\!\!
    .
\end{equation}
To derive the MoT subtraction prescription, we instead start by replacing either the $\ell^2$ propagator or the $(\ell-q)^2$ propagator with $1/q^2$.
The idea is to mimic the soft behaviour of the integral as either the $\ell^2$ propagator goes to zero or the $(\ell-q)^2$ propagator.
This procedure results in two identical triangle functions
    \begin{align}
    &\delta_{2\gamma}^{\rm MoT'} = \delta_{2\gamma} -
        \frac{2\alpha}{\pi} q^2 \Re\Bigg[p_1\cdot k_1 \nn\\
      &  \hspace{-0.5cm}\times\int  \frac{\dd \ell}{
            \ell^2
            q^2
            \Big[(\ell-k_1)^2-m^2\Big]
            \Big[(\ell+p_1)^2-M^2\Big]
        } \nn\\&\qquad\quad - (s\rightarrow u)\Bigg].
        \label{eq:motorgi}
\end{align}
Note that this is not yet the correct MoT prescription, hence {\rm MoT$'$}.
The loop integral is slightly more involved than the previous one but can still be trivially written in terms of the well-known Ellis-Zanderighi functions~\cite{Ellis:2007qk}.
In this case, we need the integral $I_6^{\rm fin}(s;m^2,M^2)$ commonly referred to as Triangle 6,
\begin{eqnarray}
    \delta_{2\gamma}^{\rm MoT'} &=& \delta_{2\gamma} +
        \frac{2 \alpha}{\pi} \Bigg\{
            (m^2 + M^2 - s) \Re\Bigg[
                I_6^{\rm fin}(s;m^2,M^2) \nn\\
                &+& \frac1{\sqrt{\lambda_s}}\ln\frac{m^2+M^2-s + \sqrt{\lambda_s}}{2mM} 
            \begin{cases} \ln\frac{m_\gamma^2}{mM} \\ \frac1{\epsilon}\Big(\frac{\mu^2}{mM}\Big)^\epsilon \end{cases}\!\!\!\!\!\!\Bigg]\nn\\
        &-& (s\rightarrow u)\Bigg\}.
\end{eqnarray}
If one chooses, $I_6^{\rm fin}$ could be expressed in terms of logarithms and dilogarithms.
The resulting expression is not overly complicated but reproducing it here serves no practical purpose.
However, in the limit of $m^2\ll s,M^2,Q^2$ the expression is fairly compact
\begin{eqnarray}
    I_6^{\rm fin}(s;m^2,M^2) &=& -\frac{y}{2M^2}\Bigg(
        -\frac13\big(i\pi+\log(-y^3)\big)^2\nn\\ &+ &2\ln(-y) \ln\frac y{1+y} - 2{\rm Li}_2(-y) \nn\\&-& \ln(-y)\ln\frac{m^2}{M^2}
    \Bigg) + \mathcal{O}(m^2)\,,
\end{eqnarray}
where we have defined $y=M^2/(s-M^2)$ and used the usual definitions of the $\zeta$ function and dilogarithm.

To arrive at the correct MoT prescription~\cite{Mo:1968cg} (cf. also~\cite{Arrington:2011dn}) we need to flip $k_1\to-k_1$ in the first term of~\eqref{eq:motorgi}, resulting in
\begin{align}
    &\delta_{2\gamma}^{\rm MoT} = \delta_{2\gamma} -
        \frac{2\alpha}{\pi} q^2 \Re\Bigg[\nn\\
      &  \hspace{-0.5cm}\phantom{+}\int \dd \ell \frac{-p_1\cdot k_1 }{
            \ell^2
            q^2
            \Big[(\ell+k_1)^2-m^2\Big]
            \Big[(\ell+p_1)^2-M^2\Big]
        } \nn\\
      &  \hspace{-0.5cm}+\int \dd \ell \frac{-p_1\cdot k_2 }{
            \ell^2
            q^2
            \Big[(\ell+k_2)^2-m^2\Big]
            \Big[(\ell+p_1)^2-M^2\Big]
        } \Bigg].
\end{align}
This operation is equivalent to setting $s\to s'=2 m^2 + 2 M^2 - s$, which spoils the crossing symmetry but eliminates the $(i\pi)^2$ from the analytic continuation of the logarithm.
The resulting expression can still be written in terms of $I_6^{\rm fin}$ but now in terms of $s'$ rather than $s$
\begin{align}
    \delta_{2\gamma}^{\rm MoT} &= \delta_{2\gamma} +
        \frac{2 \alpha}{\pi} \mathrm{Re}\Bigg [
              (s - m^2 - M^2) I_6^{\rm fin}(s';m^2,M^2) \nn\\
            &\qquad+ (u - m^2 - M^2) I_6^{\rm fin}(u;m^2,M^2)\nn\\
            &+\Bigg[
                  \frac{s-m^2-M^2}{\sqrt{\lambda_{s}}}\ln\frac{-m^2-M^2+s + \sqrt{\lambda_{s}}}{2mM}\nn\\
               &+\frac{u-m^2-M^2}{\sqrt{\lambda_{u}}}\ln\frac{ m^2+M^2-u + \sqrt{\lambda_{u}}}{2mM}
            \Bigg]\nn\\
            &\begin{cases} \ln\frac{m_\gamma^2}{mM} \\ \frac1{\epsilon}\Big(\frac{\mu^2}{mM}\Big)^\epsilon \end{cases}\!\!\!\!\!\!\Bigg\} \Bigg].
\end{align}
The symmetry of the $\lambda$ functions means that $\lambda_{s'}=\lambda_s$, which helps to simplify the result.
The difference between $\delta_{2\gamma}^{\rm MoT}$ and $\delta_{2\gamma}^{\rm MoT'}$ is
\begin{eqnarray}
    &\delta_{2\gamma}^{\rm MoT'}-\delta_{2\gamma}^{\rm MoT} = \dfrac{2\alpha}\pi (m^2+M^2-s)\Re\Bigg[\nn\\
            &+I_6^{\rm fin}(s';m^2,M^2) + I_6^{\rm fin}(s;m^2,M^2)\Bigg].
\end{eqnarray}
Crucially, the poles agree after taking the real part so that the IR cancellation in $\delta_{2\gamma}$ is unaffected.
